\bindcontent{
\section*{Question 1}

\begin{enumerate}[label=(\alph*)]
    \item 
    \begin{enumerate}[label=\roman*.]
        \item \textbf{Population:} All UCF students
        \item \textbf{Sample:} The 200 UCF students who volunteered
        \item \textbf{Parameter:} $\mu = 65$ (The true average stress level of all UCF students)
        \item \textbf{Statistic:} $\bar{x}_{before} = 66$ and $\bar{x}_{after} = 54$ (The average stress levels of the sample)
    \end{enumerate}

    \item 
    \begin{enumerate}[label=\roman*.]
        \item Quantitative, Discrete
        \item Qualitative
        \item Quantitative, Continuous
        \item Qualitative
    \end{enumerate}

    \item 
    \begin{enumerate}[label=\roman*.]
        \item Interval
        \item Nominal
        \item Ratio
    \end{enumerate}

    \item \textbf{Designed experiment.} The researcher is actively applying a treatment (the wellness program) to the subjects to measure its effect on their stress levels, rather than just observing them without interference.

    \item \textbf{No.} Because the participants were volunteers (not randomly assigned to a treatment and control group), there may be confounding variables (such as motivation to reduce stress) that caused the reduction. Without randomization and a control group, causation cannot be definitively established.

    \item 
    \begin{enumerate}[label=\roman*.]
        \item Simple random sampling
        \item It removes selection bias and ensures the sample is more representative of the entire population, unlike a volunteer sample which is prone to volunteer bias.
    \end{enumerate}
\end{enumerate}
}
